\documentclass[11pt,a4paper]{article}
\usepackage{fontspec,amsmath,amssymb,amsthm,microtype}
\setmainfont{lmroman10-regular.otf}[BoldFont=lmroman10-bold.otf,ItalicFont=lmroman10-italic.otf,BoldItalicFont=lmroman10-bolditalic.otf]
\usepackage[margin=25mm]{geometry}
\usepackage[unicode,hidelinks]{hyperref}
\hypersetup{pdfauthor={LangevinLean Project}}
\setlength{\emergencystretch}{3em}
\setlength{\parskip}{0.4em}
\setlength{\parindent}{0pt}
\allowdisplaybreaks[2]
\raggedbottom
\title{\texorpdfstring{The exact third-order LMC theorem for the actual Gaussian algorithm}{The exact third-order LMC theorem for the actual Gaussian algorithm}}
\author{LangevinLean Project}
\date{}
\begin{document}\maketitle
Draft mathematical exposition; these notes may contain mistakes.


Laplacian theorem (source th:3) follows from the actual one-step coupling and the solved quadratic recurrence, with precisely three continuous derivatives and the paper’s Laplacian-gradient condition. No Hessian Lipschitz or fourth-derivative condition is added.



The Laplacian theorem (source th:3, rendered Theorem 2) assumes \(f\in C^3\), \(m\)-strong convexity, \(L\)-Lipschitz gradient, \(0<h<1/L\), and exactly the trace-type third-order condition


\[\|\nabla\Delta f(x)\|^2\le M_\Delta^2p\qquad(x\in\mathbb R^p).\]


The actual constant-step Gaussian LMC marginals then satisfy, also for arbitrary initial laws in extended transport,


\[W_2(\nu_k,\pi)\le(1-mh)^kW_2(\nu_0,\pi)+\bigl(2L\sqrt{M_{\rm av}}+M_\Delta\bigr)\frac{h\sqrt p}{m}.\]
\section{\texorpdfstring{Pass the quadratic-plus-linear estimate to the transport infimum}{Pass the quadratic-plus-linear estimate to the transport infimum}}


For each positive \(\varepsilon\), the definition of Wasserstein distance supplies a finite initial coupling with cost at most \(W_{2}+\varepsilon\). Nonnegative quadratic and linear coefficients allow both powers of that cost to be replaced by \(W_{2}+\varepsilon\). The resulting polynomial upper bound is continuous as \(\varepsilon\) decreases to zero. Passing to this limit proves the exact quadratic-plus-linear estimate in the actual Wasserstein distance without assuming that any optimal coupling exists.



Earlier results used here: \cite[Section~3.17.1]{LangevinLean2026ULASC}; \cite[Section~3.17.2]{LangevinLean2026ULASC}; \cite[Section~3.17.3]{LangevinLean2026ULASC}; \cite[Section~3.17.4]{LangevinLean2026ULASC}; \cite[Section~3.1.1]{LangevinLean2026IPLA}.

\section{\texorpdfstring{Instantiate the error recurrence with actual marginal laws}{Instantiate the error recurrence with actual marginal laws}}


Finite initial transport propagates along the actual constant-step Gaussian LMC laws. At every step, use the canonical coupling estimate and the quadratic-plus-linear infimum theorem to obtain the scalar recurrence for their actual Wasserstein errors. Strong convexity and coordinate smoothness give \(m\le M_{\mathrm{av}}\le L\). Split \(\sqrt{M_{\mathrm{av}}p}\) into \(\sqrt{M_{\mathrm{av}}}\sqrt p\) and apply the proved recurrence with \(A=\sqrt{M_{\mathrm{av}}}\), \(\delta =M_\Delta\) and \(S=\sqrt p\). Its invariant floor is exactly \((2L\sqrt{M_{\mathrm{av}}}+M_\Delta )h\sqrt p/m\), while its transient is \((1-mh)^{k}\) times the true initial distance. The constant-schedule law identity converts the recursive marginals to iteration of the actual Gaussian transition.



Earlier results used here: \cite[Section~3.17.3]{LangevinLean2026ULASC}; \cite[Section~3.17.4]{LangevinLean2026ULASC}; \cite[Section~3.26.3]{LangevinLean2026ULASC}; \cite[Section~3.36.7]{LangevinLean2026ULASC}.

\section{\texorpdfstring{Include arbitrary initial probability laws exactly}{Include arbitrary initial probability laws exactly}}


The manuscript uses the strict step range \(hL<1\). Together with \(m\le L\) this makes \(1-mh\) strictly positive, so its \(k\)th power is positive for every \(k\). An infinite initial extended Wasserstein distance therefore gives an infinite right side, proving that case immediately. For finite initial distance, the actual-law theorem and derived finite output transport convert back to the extended nonnegative inequality. Thus the final statement imposes no independent initial second-moment assumption.



Earlier results used here: \cite[Section~3.17.3]{LangevinLean2026ULASC}; \cite[Section~3.17.4]{LangevinLean2026ULASC}; \cite[Section~3.26.3]{LangevinLean2026ULASC}; \cite[Section~3.36.7]{LangevinLean2026ULASC}.

\begin{thebibliography}{ULA-LSI}
\bibitem[ULA-SC]{LangevinLean2026ULASC}
LangevinLean Project. \emph{Unadjusted Langevin under Strong Convexity}. LL-01, 2026. Draft report, content revision c58ab5c70c3a. \url{https://odakyildiz.github.io/LangevinLean/docs/pdf/UnadjustedLangevin/FullProof.pdf}.
\bibitem[IPLA-SC]{LangevinLean2026IPLA}
LangevinLean Project. \emph{Interacting Particle Langevin under Strong Convexity}. LL-03, 2026. Draft report, content revision ce6396dff331. \url{https://odakyildiz.github.io/LangevinLean/docs/pdf/IPLA/FullProof.pdf}.
\end{thebibliography}

\end{document}
